
We began expecting a precision hierarchy --- that a more formally rigorous verifier would describe a defect more precisely and so produce better repair. Four categories, one backbone, 421 problems, and a Spearman correlation of -1.0 (n=4 categories; see §5.4 for caveats) later, that expectation is not weakly supported but exactly inverted. The verifier that says the most fixes the least, and the verifier that costs the least returns the most correctness per second.

Holding GPT-4o-mini, the problem set, the prompt template, and a three-iteration budget fixed while varying only the feedback category, we found: the four categories produce feedback volumes separated by four orders of magnitude (Kruskal-Wallis $\epsilon ^2$ = 0.88), which makes the comparison meaningful; iteration-1 repair success falls monotonically as feedback volume rises, from Z3's 7.94\% at the top down to Pyright's 4.76\%, giving $\rho$ = -1.0; normalizing correctness by wall-clock overhead makes static analysis the efficiency winner at an estimated 6.64 ratio against execution monitoring's estimated 0.41, despite execution monitoring delivering twice the absolute pass@1 improvement; and SMT-guided repair is not reachable at this model tier at all, with usable Z3 constraints extractable for 6.3\% of problems against the 40\% assumed.

Our account of the inversion is that repair is bounded by the model's ability to extract one actionable edit, not by the feedback's information content. A short traceback names the fix; a 24,000-character diagnostic tree contains it somewhere.

Two of these results came from experiments that failed pre-registered gates. We report them as the contribution because a refuted prediction that reverses a field-wide assumption is worth more than a confirmed one that restates it.

The most urgent future work separates the two explanations for the inverse correlation. Comparing raw Pyright JSON against a one-sentence model-generated summary of the same diagnostics, on matched problems, would settle whether the problem is presentation or truncation: if the summary matches execution monitoring's repair rate while the raw JSON does not, the fix is a summarization step costing one extra call. Running the three-category comparison on GPT-4o and Claude 3.5 Sonnet would test whether the mechanism generalizes across model tiers --- a sign reversal would invert the deployment advice rather than merely qualify it. Beyond that, the bug distribution (65.9\% logic, 24.6\% type at 82\% classifier agreement) suggests a routing strategy: static analysis on predicted type errors, execution monitoring on logic errors, which would cost roughly the same as the repair experiment to check.

The verify stage of a repair loop is usually designed around what the tool can prove. Our results suggest it should be designed around what the model can act on. Those turn out to be different objectives, and on current models they point in opposite directions.

